\section{Memory：从上下文缓存到可持续知识底座}

\subsection{记忆分层设计}
课程将 Memory 视为 Agent 稳定性的决定性因素。没有分层记忆，长任务会退化成 ``每轮从零开始''。本讲建议至少区分：working memory、episodic memory、semantic memory、procedural memory。

\begin{figure}[H]
\centering
\includegraphics[width=0.84\textwidth]{slide-04-memory-design.jpg}
\caption{视频约 00:52:40 讲义页：记忆模块不是单一向量库，而是层次化读写系统}
\end{figure}

\begin{table}[H]
\centering
\begin{tabular}{p{2.8cm}p{3.6cm}p{4.8cm}}
\toprule
\textbf{记忆层} & \textbf{典型内容} & \textbf{工程重点} \\
\midrule
Working memory & 当前步骤状态、临时变量 & 快速更新、低延迟、可裁剪 \\
Episodic memory & 任务轨迹、失败记录、决策分支 & 支持复盘与策略迁移 \\
Semantic memory & 稳定事实、领域知识、工具文档 & 检索精度与版本一致性 \\
Procedural memory & 可复用工作流、策略模板、工具调用范式 & 抽象粒度与触发条件设计 \\
\bottomrule
\end{tabular}
\caption{四层记忆在 Agent 系统中的分工}
\end{table}

\begin{importantbox}{Memory 的本质是可控读写，不是单向存储}
记忆系统必须回答四个问题：写什么、何时写、如何读、何时忘。若只做 ``全量写入''，很快会出现噪声累积和检索错配，最终拖垮推理与规划质量。
\end{importantbox}

\begin{knowledgebox}{Memory 操作的最小 API}
\begin{itemize}
    \item `write(event, importance, scope)`：写入时附带重要性与作用域；
    \item `retrieve(query, budget, freshness)`：检索时显式预算和时效约束；
    \item `compress(window)`：按任务阶段进行摘要与压缩；
    \item `evict(policy)`：执行过期淘汰与冲突合并。
\end{itemize}
\end{knowledgebox}

\begin{warningbox}{记忆污染是长任务失败的隐蔽主因}
若历史错误结论被反复检索，Agent 会稳定地做错事。应通过来源置信度、冲突检测和回滚机制降低污染传播。
\end{warningbox}

\subsection{本章小结}
Memory 决定了 Agent 能否跨轮次、跨任务累积有效能力。分层、可控、可回滚是高质量记忆系统的三个底线。

